\kbvchapter{ಪಾಠ ೧೨}{ಕ್ರಿಯಾಪದ-ಲಿಂಗ, ವಚನ, ಕಾಲ}

\noindent\hspace*{2.2em}ಕೆಳಗೆ ಕೊಟ್ಟಿರುವ ವಾಕ್ಯಗಳನ್ನು ನೋಡಿರಿ.\par

\indent - ಹುಡುಗನು ರೊಟ್ಟಿ ತಿನ್ನುತ್ತಿದ್ದಾನೆ. (ಪು. ಏ.)\par

\indent - ಹುಡುಗರು ರೊಟ್ಟಿ ತಿನ್ನುತ್ತಿದ್ದಾರೆ. (ಪು. ಬ.)\par

\indent - ಹುಡುಗಿಯು ರೊಟ್ಟಿ ತಿನ್ನುತ್ತಿದ್ದಾಳೆ. (ಸ್ತ್ರೀ.)\par

\indent - ಹುಡುಗಿಯರು ರೊಟ್ಟಿ ತಿನ್ನುತ್ತಿದ್ದಾರೆ. (ಸ್ತ್ರೀ. ಬ.)\par

\indent - ನಾಯಿಯು ರೊಟ್ಟಿ ತಿನ್ನುತ್ತಿದೆ. (ನ. ಏ.)\par

\indent - ನಾಯಿಗಳು ರೊಟ್ಟಿ ತಿನ್ನುತ್ತಿವೆ. (ನ. ಬ.)\par

\indent - ಮಕ್ಕಳು ರೊಟ್ಟಿ ತಿನ್ನುತ್ತಿದ್ದಾರೆ. (ನ. ಬ.)\par

\indent ಈ ವಾಕ್ಯಗಳಲ್ಲಿ “ತಿನ್ನು” ಧಾತುವಿನ ಮೇಲೆ ಲಿಂಗ ವಚನಗಳಿಗೆ ಅನುಸಾರವಾಗಿ ಬೇರೆಬೇರೆ ಆಖ್ಯಾತ ಪ್ರತ್ಯಯಗಳು ಸೇರಿ ಕ್ರಿಯಾಪದಗಳಾಗಿರುವುದನ್ನು ಕಾಣಬಹುದು.\par

\indent ಮೊದಲನೆಯ ವಾಕ್ಯದಲ್ಲಿ “ಹುಡುಗನು” ಎಂಬ ಕರ್ತೃಪದವು ಪುಲ್ಲಿಂಗ ಏಕವಚನದಲ್ಲಿರುವುದು ; ಆದುದರಿಂದ ಕ್ರಿಯಾಪದವೂ ಪುಲ್ಲಿಂಗ ಏಕವಚನದಲ್ಲಿರುವುದು.\par

\indent ಹೀಗೆ, ಪ್ರತಿಯೊಂದು ವಾಕ್ಯದಲ್ಲೂ ಕರ್ತೃಪದವು ಯಾವ ಲಿಂಗದಲ್ಲಿರುವುದೋ ಕ್ರಿಯಾಪದವೂ ಅದೇ ಲಿಂಗದಲ್ಲಿರುವುದು ; ಅಲ್ಲದೆ, ಕರ್ತೃಪದವು ಯಾವ ವಚನದಲ್ಲಿರುವುದೋ ಕ್ರಿಯಾಪದವೂ ಅದೇ ವಚನದಲ್ಲಿರುವುದು. ಕ್ರಿಯಾಪದವು ಕರ್ತೃವಿನ ಕೆಲಸವನ್ನು ತಿಳಿಸುವುದರಿಂದ, ಕರ್ತೃ ಹಾಗೂ ಕ್ರಿಯಾಪದಗಳಲ್ಲಿ ಲಿಂಗ ವಚನಗಳು ಏಕರೀತಿಯಾಗಿರಬೇಕಾದ್ದು ಅತ್ಯಾವಶ್ಯಕ.\par

\indent ಕಾಲ ಸೂಚಕ ಪ್ರತ್ಯಯಗಳು\par

\vspace{4pt plus 4pt minus 1pt}\noindent{\kbvsecfont\fontsize{11pt}{14.5pt}\selectfont\color{kbvaccentdark} (1)\hspace{0.45em}ವರ್ತಮಾನ ಕಾಲದ ಕ್ರಿಯಾಪದ ರೂಪಗಳು :}\hspace{0.5em}ಕ್ರಿಯೆ ಈಗ ನಡೆಯುತ್ತಿರುವುದಾಗಿ ತಿಳಿಸುವ ಕ್ರಿಯಾಪದವು. ವರ್ತಮಾನ ಕಾಲದ ಕ್ರಿಯಾಪದವೆನಿಸುವುದು.\par
